\label{night:P3K.96}
\StatementRecord{night:P3K.96}{proved}{}{\cref{night:P3F.35}, \cref{night:P3F.41}, \cref{night:P3F.46}, \cref{night:P3K.94}}
Retain the actual source
\[
 y^2=x^3+r x^2+t x+t^2,\qquad G=(0,t),\qquad
 M_a=\mathcal O(aG-O),\qquad Q=bG,
\]
on a sufficiently small disc about \(r=0\). Choose the analytic branch
\(\rho^2=1-3r\), \(\rho(0)=-1\), and the selected whole
\(I_1\) zero \(t_c=(1-3\rho^2+2\rho^3)/27\). For either of the
complete cap progressions
\[
\begin{array}{c|c|c|c}
(a,b)&m&\text{root fibre}&\text{range}\\\hline
(6,1)&3l+2&IV&l\ge0\\
(2,3)&3l&\text{smooth elliptic}&l\ge1
\end{array}
\]
use the exact Cartier suspension
\(B_{\rm sel}=(t-t_c)B_{a,b}/t^3\) of the proof, either primitive
norm \(n=\pm1\), every prescribed character \(\zeta_m^d\), and
the entire specified reference annular map. Retain a fixed sufficiently
small nonzero scalar for the chosen fan, and no extra logarithmic disc
at the colliding fibres. For every prescribed lattice-basis fan at the
selected cusp, these actual maps extend to a smooth proper oriented
sphere family over the full source disc, with identity parameter
closing. Its positive boundary return is \(0\in\Theta_7\).
The literal marked-cover map \(D=w^2-3\),
\(r=(3-2w^2)/D^2\), \((X,Y,T)=(D^2x,D^3y,D^3t)\)
carries the same conclusion to both points \(w^2=3/2\), with the
full scalar, line, cap and fan marking retained.
\end{theorem}
\begin{proof}
\begin{enumerate}[label=\arabic*.]
\item \par

\medskip\noindent\textbf{EX.1. Full point-line multiplier and its finite vertical divisor}

\par\nobreak

\item Use the same actual source
\[
\begin{gathered}
 y^2=x^3+r x^2+t x+t^2,\quad G=(0,t),\quad s=\tfrac14-r,\\
 U=y-t-x/2,\quad f_2=\frac U{x-s},\quad T(p)=p-G.
\end{gathered}
\]
The typed point-to-tensor arrow is
\(r_2=f_2^{-1}:M_2\longrightarrow M_1^{\otimes2}\).
It gives the literal single-step lift
\[
 B_{2,1}(p)=f_2(p)^2\frac{f_2(Tp)}{f_2(p)}
          =f_2(p)f_2(Tp).
\]
Composing the three actual source/line arrows, without cancelling
their source or target Cartier factors separately, gives
\[
 B_{2,3}(p)=f_2(p)f_2(T^3p)f_2(Tp)^2f_2(T^2p)^2.
 \tag{EX.1}
\]
Its source map is exactly \(T^3=p-3G\).

\item At a generic I2 identity component \(C_0\), the functions \(U\)
and \(x-s\) are generic units, so \(\operatorname{ord}_{C_0}f_2=0\).
At the other component use the actual chart \(x=t\alpha,y=t\beta\):
\[
 f_2=t\frac{\beta-1-\alpha/2}{-s+t\alpha}.
\]
The displayed numerator and denominator are generic units and
\(s\ne0\), hence \(\operatorname{ord}_{C_1}f_2=1\).
The exact inverse translation of K94 EU.0 exchanges the two
components. Therefore the two valuations of \(B_{2,1}\) are
both one, and those of EX.1 are both three. The whole finite
vertical divisor at zero is exactly \(3F_0\).

\item At any other finite irreducible fibre, neither \(U\) nor \(x-s\)
contains that fibre: substituting \(U=0\) in the source gives
\(x^2(x-s)=0\), and the fibre is not one of these point divisors.
Every power of \(T\) preserves the generic component DVR and its
base uniformizer. Thus each factor in EX.1 is a DVR unit there.
At infinity \(3G\) fixes each component, so the Hom multidegrees
vanish. The primitive affine component kernel gives a whole fibre
there; its total Hom degree is \(2\cdot3\cdot(1/6)=1\).
Consequently its remaining vertical divisor is \(-2F_\infty\).
The exact selected-zero map is therefore
\[
 B^{(2,3)}_{\rm sel}=\frac{t-t_c(r)}{t^3}B_{2,3}.
 \tag{EX.2}
\]
The actual simultaneous A1 source, the maps \(T,T^{-1}\), and
the section \(2G=(s,-t-s/2)\) extend on the r-disc of K94.
EX.2 and its inverse have zero order at the parameter divisor,
as seen on an ordinary source section avoiding the horizontal
divisors. The only remaining locus is the codimension-two central
curve. The same height-one-prime argument as K94 extends this
actual Hom map and its inverse across it, with its prescribed
selected zero retained. No parameter scalar is changed.

\par

\medskip\noindent\textbf{EX.2. The literal good root and its central multiplier}

\par\nobreak

\item For \(m=3l\), \(l\ge1\), the actual good root is
\[
 q_\infty=z^{3l},\quad x_0=z^{4l}X,\quad y_0=z^{6l}Y,
 \qquad Y^2=X^3+r z^{2l}X^2+z^lX+1.
 \tag{EX.3}
\]
It has \(G'=(0,1)\) and
\[
 2G'=(s z^{2l},-1-(s/2)z^{3l}).
\]
The deck acts by \((z,X,Y)\mapsto
(\zeta_m z,\zeta_m^{-4l}X,Y)\), retaining the two marked
three-torsion limits and the actual line \(M_2\).
The original coordinates are \(x=z^{-2l}X,y=z^{-3l}Y,t=z^{-3l}\).
Hence each tangent-over-vertical factor has exact weight
\(f_{2,\rm old}=z^{-l}f'_2\). All six factors in EX.1 give
\[
 B_{2,3,\rm old}=z^{-6l}B'_{2,3},\qquad
 B^{(2,3)}_{\rm sel}=h B'_{2,3},\qquad
 h=1-t_c(r)z^{3l}.
 \tag{EX.4}
\]
This is the full coefficient scaling in the actual rational
source and target Cartier frames.

\item On the central good elliptic curve \(Y^2=X^3+1\), put
\(f=(Y-1)/X\), \(T'=p-G'\). Since \(G'\) has order three,
\(T'^3=1\), and the product of the three translated divisors
of \(f\) is zero. Its norm is therefore a nonzero constant.
The actual orbit of \((2,3)\) is
\((2,3),(2,-3),(-1,0)\); the three values of \(f\) are
\(1,-2,1\). Thus its norm is exactly \(-2\), and
\[
 B'_{2,3}|_{z=0}
   =[f(p)f(T'p)f(T'^2p)]^2=4.
 \tag{EX.5}
\]
The central selected coefficient in EX.4 is consequently \(4h\),
with its full value retained.
The same constant follows from the entire actual central map:
\[
 T'(X,Y)=\left(\frac{2X}{Y-1},-\frac{Y+3}{Y-1}\right),\qquad
 T'^2(X,Y)=\left(-\frac{2X}{Y+1},\frac{Y-3}{Y+1}\right).
\]
These formulae have \(T'^3=1\), and the three factor values as
rational functions are \((Y-1)/X,-(Y+1)/X,2/X\). Their product
is \(-2\) by the exact curve equation \(Y^2-1=X^3\).

\item The section \(Q'=3G'\) is in the identity formal group, with
\(\tau(Q')=-s z^{2l}/2+O(z^{5l})\). Its small analytic
\(m\)-division exists and is deck invariant by uniqueness.
The full good-root point-line correspondence, its invariant toric
correction and norm lift are constructed in EX.5. EX.6 carries
the prescribed finite character and entire annular reference
through the actual disc attachment.

\par

\medskip\noindent\textbf{EX.3. Marked-cover coefficient check}

\par\nobreak

\item Under the source map \(X=D^2x,Y=D^3y,T=D^3t\), every factor
\(f_2\) has weight \(D^{-1}\), so EX.1 has weight \(D^{-6}\).
The selected-height factor has weight \(D^6\), exactly as in
K94 EU.14. Therefore EX.2 is carried with scalar unchanged.
This actual Miller/Cartier coefficient equality carries the
constructed disc family to the marked w-cover without changing
its scalar or selected toroidal fan.

\par

\medskip\noindent\textbf{EX.4. Every reduced root: actual middle charts and full line graph}

\par\nobreak

\item The initial local stage is preserved in \path{records/attempts/e6_unselected_iii_extension/initial_local.md}. The following maps complete both cap progressions.
Keep the same source-parameter disc, selected whole zero, fixed
small scalar and full reference annular map as in K94. The two
named data are
\[
 (a,b;m)=(6,1;3l+2),\quad l\ge0,
 \qquad (a,b;m)=(2,3;3l),\quad l\ge1,
\]
with either signed sphere norm \(n=\pm1\) and the retained finite
character \(\zeta_m^d\). There are no logarithmic discs at the
colliding fibres. The following maps construct the full marked disc family.

\item For the reduced branch put \(m=3l+2,k=2l+1\). The literal root is
\[
 q_\infty=z^m,\quad x_0=z^{2k}X,\quad y_0=z^{3k}Y,
\]
\[
 Y^2=X^3+r z^{2l+2}X^2+z^{l+2}X+z^2,\qquad G'=(0,z).
 \tag{EX.6}
\]
Write \(b_0=z^l,c_0=r z^{2l}\) and
\(c=\sqrt{1+b_0X+c_0X^2}\). The exact coordinates
\(U=Y-zc,V=Y+zc\) satisfy \(UV=X^3\). The equation
\(2zc=V-U\) has a unique analytic inverse for \(z\) near the
origin, since its derivative there is two. Thus the same three
minimal A2 resolution charts give a smooth family over the entire
source disc, retaining the actual function \(q_\infty=z^m\).

\item In the middle chart \(X=ab,U=a^2b,V=ab^2\), set
\[
\begin{aligned}
 K&=1+\frac{b(b-a)(b_0+c_0ab)}{2c(c+1)},\\
 F&=1-\frac{ab(b_0+c_0ab)(K+1)}{2c(c+1)}
           -\frac{c_0ab^2(a-b)}{4c^2},\\
 X'&=a(a-b)F,\qquad c'=\sqrt{1+b_0X'+c_0X'^2},\\
 H&=K+\frac{b(a-b)(b_0+c_0X')}{2c(1+c')}.
\end{aligned}
\]
The actual \(+G'\) chord has \(\kappa=(Y-z)/X=aK\),
\(X'=\kappa^2-c_0z^2-X\), \(Y'=-z-\kappa X'\).
Its complete toric chart arrow is
\[
                  a'=(b-a)F/H,\qquad b'=-aH.
 \tag{EX.7}
\]
As in K94, \(V'=Y'+zc'=-aX'H\) proves the formula, and the
analytic inverse of \(2zc=V-U\) proves that it retains exactly
the same \(z\). All displayed denominators are units and
\(K=F=H=1\) at the triple point. The inverse is obtained from
the actual \(-G'\) chord, equivalently by \(Y\mapsto-Y\).
The smooth group charts at the other central points complete
this actual automorphism.

\item The full deck is
\[
 D(z,X,Y)=(\zeta_mz,\zeta_m^{-2k}X,\zeta_m^{-3k}Y).
\]
Since \(l-2k=-m\), \(c\) is deck invariant; \(U,V\) have
weight \(\zeta_m\). In the middle chart both \(a,b\) have
weight \(\zeta_m^{l+1}\). The identities
\(3(l+1)=m+1\) and \(\gcd(l+1,m)=1\) verify the actual
order and inverse-base relation. Every term of \(K,F,H\) has
weight zero. Formula EX.7 therefore commutes with this whole
deck, including its parameter coordinate.

\item The invariant vertical field has an explicit regular lift there.
Put \(J_0=b_0+2c_0ab\). The usual Weierstrass vector field gives
\[
 W=a\left(2b-a-\frac{zJ_0}{c}\right)\partial_a
   +b\left(2a-b+\frac{zJ_0}{c}\right)\partial_b.
 \tag{EX.8}
\]
This follows from
\(W(U)=3X^2-zJ_0U/c\), \(W(V)=3X^2+zJ_0V/c\).
It is holomorphic and vanishes at the triple point. Its flow
agrees with the identity-component group flow on the smooth
locus. Thus the exact division section
\[
 R=2G'-[m]^{-1}((2m-1)G'),\qquad mR=G'
 \tag{EX.9}
\]
extends as an actual source translation on the full resolved
root. Here \(2m-1=3(2l+1)\), so the argument of the inverse is
in the small identity formal group; the inverse has derivative
\(1/m\). The central component of \(R\) is two. Both EX.9
and the resulting translation commute with \(D\), by uniqueness
of that formal inverse and the displayed deck arrows.

\item For the line, retain the exact source coordinate
\(\tau_0=-x_0/y_0\) and root coordinate \(\tau'=-X/Y\).
Their relation and the full section equation are
\[
 \tau'=z^k\tau_0,\qquad
 \tau'(6G')=z^k L(s,z^m),\qquad L(s,0)=-s\ne0.
 \tag{EX.10}
\]
Use the normalized closure of the graph of the displayed generic
map to the actual minimal regular source. Its two projections
are the morphisms denoted \(\beta\) to the root and \(\pi\)
to the source. Near the marked sections this is the graph of
\(u=\tau'/z^k\), constructed by the actual ideal
\((\tau',z^k)\). The root section equations are
\(\tau'\) and \(\tau'-z^kL(s,z^m)\). The normalized graph adds the ray \((k,1)\). Its cone with
\((1,0)\) is regular. Subdivide the other cone by the literal
rays \((j,1)\), \(1\le j\le k\), followed by \((0,1)\).
Adjacent ordered determinants are
\(j\cdot1-(j-1)\cdot1=1\); these are exactly the exceptional
rays of this chosen regular graph. On the ray \((j,1)\),
\(\operatorname{ord}(\tau')=j\) and
\(\operatorname{ord}(\tau'-z^kL)=\min(j,k)=j\).
At the equality ray \(j=k\),
the graph coordinate \(u\) retains their distinct directions
\(0\) and \(L(s,q_\infty)\); its two strict section equations
are carried separately. The literal coefficient identity is
\[
 \frac{\tau'}{\tau'-z^kL(s,z^m)}
       =\frac{u}{u-L(s,q_\infty)}.
 \tag{EX.11}
\]
On the other graph chart, \(\tau'v=z^k\), the same actual
coefficient is \(1/(1-vL(s,z^m))\). It is a unit at \(v=0\).
The two equations agree by \(uv=1\); the remaining cyclic
singularity of \(\tau'v=z^k\) is resolved by its ordinary
monomial subdivision, and both coefficient arrows pull back
unchanged to every resulting chart. This explicitly retains
the direction at equal valuation as well as both outer charts.
Both marked divisors are absent from the other source/root
quotient centres. Resolving those remaining cyclic charts by
their invariant monomial ideals does not add a coefficient to
either divisor. Thus on the full common graph, and on each of
its such refinements,
\[
                \beta^*N_6=\pi^*M_6.
 \tag{EX.12}
\]
This equality is an equality of Cartier line maps, not only of
generic fibre classes. The source and root translations preserve the graph by the
following actual square. Since the root translation commutes with
\(D\), it and its inverse descend to the finite quotient. An automorphism
of this quotient lifts to its minimal resolution: the local cyclic
resolution is given by the invariant monomial ideals already
specified, and their exceptional valuations are preserved. The
subsequent contraction to the unique minimal regular source is
also preserved, since the product carries each contractible
vertical curve and its intersection data to that of its image.
Consequently there are actual mutually inverse source maps over
\(q_\infty\), extending the descended translation. The product of the
source and root maps carries the generic graph to itself, as does
its inverse. Taking its closure and normalization therefore gives
both inverse graph maps. The marked-section graph ideals and the
displayed cyclic subdivisions carry these maps, retaining the
actual Cartier coefficient square on every chart.

\item The exact root multiplier is also uniform in \(l\). The original
coordinates satisfy \(x=z^{-2(l+1)}X,y=z^{-3(l+1)}Y,t=z^{-m}\).
The Miller weights therefore give
\[
 B_{\rm sel}'=h z^{-2}B'_6,\qquad h=1-t_c(r)z^m.
 \tag{EX.13}
\]
For a generic root identity-component point \(X=w^2,Y=w^3+O(z^2)\),
the full three-point factor is
\[
 f'_3=\frac{s z^k X(Y-z-z^{l+1}X/2)}
              {s z^k(Y+z+z^{l+1}X/2)-(s z^m+2)X}.
\]
The actual inverse chord gives
\[
 f'_2\to w,\quad f'_3(p)/z^k\to-s w^3/2,\quad
 f'_3(Tp)/z^{k+1}\to s,
\]
\[
 z^k f'_{2,3G'}(p),\ z^k f'_{2,3G'}(Tp)\to-4/s.
\]
Consequently \(B'_6/z^2\to4\). The full divisor calculation
and EX.12 give order two on every reduced central component;
normality then extends EX.13 and its inverse as actual units.
Their central coefficient is \(4h\), with no omitted base
factor. The case \(l=0\) is the actual quadratic map of K94.

\par

\medskip\noindent\textbf{EX.5. The full good-root Cartier correction and actual translation lift}

\par\nobreak

\item For the good branch EX.3 write \(u=z^l\). Its quotient to the
source factors through the actual minimal cubic root
\[
 Y^2=X^3+r u^2X^2+uX+1,\qquad q_\infty=u^3.
\]
The \(\mu_l\) kernel acts only on \(z\), and the remaining
\(\mu_3\) action has weights \((2,1)\) at each of its three
fixed sections. At \(2G'\) one can use
\(w=X-su^2\); at \(O\) use \(w=-X/Y\). They have exactly
weight two, with \(u\) of weight one. Their quotient lattice is
\(\mathbb Z^2+\mathbb Z(2,1)/3\), with the two primitive
resolution rays \((2/3,1/3),(1/3,2/3)\).

\item Here is the actual strict-section Cartier calculation. On that
resolved quotient its support function is
\(h(v_w,v_u)=\max(v_w-2v_u,0)\): it is \(v_w-2v_u\)
on the cone adjacent to the section axis and zero on the other
cones. Pulling the section to the common root fan has coefficient
\(v_w\). Their difference is therefore
\(\min(v_w,2v_u)\). In the full \(z\)-root it is
\(\min(v_w,2l v_z)\). If \(E_P\) and \(E_O\) denote
these explicit sums at the two marked sections, the full equality is
\[
          \beta^*N_2=\pi^*M_2+E_P-E_O.
 \tag{EX.14}
\]
The coefficient arrow is the quotient of the actual monomial
Cartier generators on each cone, so it includes their units.
The same projective charts and monomials work holomorphically
over the whole \(r\)-disc; their marked centres are the actual
sections just displayed.

\item The section \(Q'=3G'\) has
\(\tau(Q')=-s u^2/2+O(u^5)\). Its selected small
\(m\)-division \(R=[m]^{-1}Q'\) has
\(\tau(R)=-s u^2/(2m)+O(u^5)\), and is deck invariant.
At each marked fixed point its translation has the exact form
\(w\mapsto w\varepsilon(w,u,r)+u^2 a(w,u,r)\), with
\(\varepsilon\) an analytic unit. It preserves the actual
ideals \((w,u^j)\) and \((u,w^j)\) for \(j=1,2\), as
does its inverse. Their invariant product and normalized blowup
give the common toric refinement above. Thus translation by
\(R\), as well as by \(Q'\), lifts to that graph and preserves
both corrections in EX.14. This proves the full equivariance of
the point-line correspondence. It retains the translated
correction through its actual canonical divisor section.

\item Both branches admit a particularly explicit starting line lift.
For \(N=\mathcal O(P-O)\), define the actual chord-over-vertical
function
\[
 g_{P,R}(p)=\frac{\ell_{P,R}(p)}{x(p)-x(P+R)}.
 \tag{EX.15}
\]
Its divisor is \((P)+(R)-(P+R)-(O)=D_P-t_{-R}^*D_P\),
so it is the coefficient of a typed line map covering
\(p\mapsto p-R\). The points used here are distinct on a
sufficiently small punctured root disc. The line functions are
the literal chord functions; no leading coefficient is removed.
On good roots take \(P=2G'\) and
\(F=\tau(R)g_{P,R}\). On reduced roots take \(P=6G'\) and
\(F=\tau(P)g_{P,R}\). The formal coordinate factors have
orders \(2l\) and \(k\), respectively, with nonzero leading
coefficients \(-s/(2m)\) and \(-s\).

\item To check the full central normalization, a point tending to
\(O\) has \(x=\tau^{-2}+O(1)\),
\(y=-\tau^{-3}+O(\tau^{-1})\). The actual chord slope is
\(-\tau^{-1}+O(1)\). At a generic test point, multiplying
the chord by \(\tau\) tends to its vertical denominator.
Thus the displayed \(F\) has central coefficient exactly one.
In the reduced case the Hom multidegrees are zero, so the
primitive central component kernel makes the chord\textquotesingle s vertical
order a whole-fibre order; the single generic-component limit
fixes it on every component. In the good case there is one
smooth component. Horizontal divisors are exactly those in
EX.15. The height-one-prime argument therefore extends \(F\)
and its inverse through the whole root family, retaining its
full normalizing factor.

\item This particular lift commutes with the deck literally. The
root Weierstrass coordinate weights are \((-2k,-3k)\), with
\(k=2l\) in the good case and \(k=2l+1\) in the reduced case.
Consequently the actual chord-over-vertical function has weight
\(\zeta_m^{-k}\), whereas each of \(\tau(P),\tau(R)\)
has weight \(\zeta_m^k\). Their displayed product has weight
one. These are identities of rational coefficient maps, so
their extending maps commute on every root chart.

\item The ratio \(\beta=B'_{{\rm sel}}/F^m\) covers the identity
on the same compact connected source fibre. It is therefore
an invertible base function, with central value four. It is
deck invariant by the preceding exact arrows. Its analytic
\(m\)-th root, fixed by the complete reference lift, is deck
invariant as well: the ratio of its transform to itself is a
constant \(m\)-th root of unity, equal to one on the central
disc. Define
\[
 \widetilde F=\beta^{1/m}F,\qquad
 A_0=\zeta_m^d\widetilde F^{\,n}D.
\]
Then, as literal source/Cartier maps,
\[
 \widetilde F^m=B'_{{\rm sel}},\qquad
 A_0^m=(B'_{{\rm sel}})^n,\qquad
 [A_0,B'_{{\rm sel}}]=1.
 \tag{EX.16}
\]
The reference finite character is still \(\zeta_m^d\).
The normalized source/tensor factors in EX.1 and the full
correction EX.14 are carried through this equality. No generic
Picard-class comparison replaces their coefficient maps.

\par

\medskip\noindent\textbf{EX.6. General-order annular square, actual disc closing and return}

\par\nobreak

\item Fix \(R_\lambda\ne0\) with \(R_\lambda^m=\lambda\),
independent of the source parameter. Put
\[
 B_\lambda=\lambda B'_{{\rm sel}},\qquad
 A_\lambda=R_\lambda^n A_0.
\]
On the two fixed annulus sectors use
\[
 s_\alpha=\frac{\log_\alpha q_\infty}{2\pi i},\quad
 z_\alpha=\exp(2\pi i s_\alpha/m),\quad
 c_\lambda=\frac{n\log R_\lambda}{2\pi i}+\frac dm.
\]
The actual map in the source/line Lie cover is
\[
 \Phi_\alpha=\operatorname{Exp}
       (n s_\alpha f_\alpha+s_\alpha c_\lambda\delta
                   +h_\alpha)I_{z_\alpha},
 \qquad f_\alpha=\operatorname{Log}_\alpha\widetilde F.
 \tag{EX.17}
\]
The source displacement is the signed \(-n s_\alpha\) Abel
integral of the actual \(R\). The scalar is
\(R_\lambda^{n s_\alpha}\exp(2\pi i d s_\alpha/m)\),
and the root change is the actual \(D\). Thus both the full
norm and every finite character appear in this formula.

\item On an overlap, with \(s_\beta=s_\alpha+k\) and
\(f_\beta=D_*^k f_\alpha+\Pi_{\alpha\beta}\), prescribe
\[
 h_\beta-D_*^k h_\alpha
   =-n s_\alpha\Pi_{\alpha\beta}
       +j_{\alpha\beta}b_\beta+\Lambda_{\alpha\beta},
 \qquad b_\beta=\operatorname{Log}_\beta B_\lambda.
 \tag{EX.18}
\]
All period columns, suspension decks and integral coefficients
are those of the full reference map. Their functions are
transported holomorphically over the source disc. Here is the reference-zero construction in the same fixed
ordinary Lie frame. Extend the prescribed reference cochains
\(h_\alpha^*\) to \(h_\alpha^0\), retaining the full
transported period functions. Let \(E_{\alpha\beta}\) be
the residual of EX.18; it vanishes at the reference parameter.
For the fixed annulus partition put
\[
 k_\alpha=-\sum_\beta\chi_\beta E_{\alpha\beta},
 \qquad \omega=\bar\partial k_\alpha,
 \qquad v(q,r)=\frac1\pi\int
                \frac{\omega(w,r)}{q-w}\,dA(w).
\]
The compactly supported residual and fixed contours are exactly
those of K94 EU.11e. The common \(\omega\) and its Cauchy
transform depend holomorphically on the source parameter and
vanish at the reference. Then
\(h_\alpha=h_\alpha^0+k_\alpha-v\) solves EX.18,
retains every integral period and suspension column, and equals
the whole prescribed reference cochain. Subtracting the target-sheet logarithm
\(kn f_\beta+kc_\lambda\delta\) from the two sides of
EX.17 leaves exactly the prescribed period and suspension deck.
Hence the actual squares are
\[
 \Phi_\beta=A_\lambda^k B_\lambda^{j_{\alpha\beta}}
                  \Phi_\alpha,\qquad
 \Phi_\alpha B_\lambda=B_\lambda\Phi_\alpha,
 \qquad A_\lambda^m=B_\lambda^n.
 \tag{EX.19}
\]
They define the ordinary/cap gluing with its actual inverse
and its full reference value. The finite cap itself extends
by EX.3, EX.6-\/-16, without extending an annular logarithm
across its centre.

\item All these maps are single valued on the whole source disc.
Their positive boundary therefore has identity source, line,
root and annular closing; the scalar is fixed. The selected
cusp remains disjoint from the collision, and its full Tate,
Cartier and cone-character maps extend exactly as in K94 EU.5
on every prescribed lattice-basis fan. On the finite charts
their actual line multipliers have compact uniform bounds.
Choosing the original fixed scalar sufficiently small gives
strict contraction. For both norm signs \(H=A_\lambda^n\)
has \(H^m=B_\lambda\) and scalar factor \(R_\lambda\).
Its compact norm shell and the protected toric fundamental
regions give the same uniform properness, freeness and smooth
disc submersion as the explicit quotient proof of K94.

\item At the central source the two data are exactly the E6-B
entries of F41/F46, with the retained ordinary III cap of
order one and norm zero. Their determinant is \(n=\pm1\),
so the full integral maps give a sphere on the reference fan.
The actual finite fan word, common blowups, two ruling maps
and primitive Gysin/van Kampen maps of F35 carry that sphere
to each prescribed fan. Radial transport in the resulting
proper disc family gives a product with its actual identity
closing. Thus the positive unselected-III boundary return is
zero for both entire cap progressions just constructed.
The quotient, line, annular and topology maps above prove the stated return for the full marked family.

\par

\medskip\noindent\textbf{Executed certificate and retained attempts}

\par\nobreak

\item The exact program \path{code/part3/e6_unselected_iii_extension.py}
checks the full central translation orbit, generic three-step
Cartier multiplier, symbolic root exponents, universal A2 chart
identities, regular graph rays and section coefficients, deck and
line weights, both signed norms and the target-sheet overlap
polynomial. Its current output records 5,241 exact identities
and finite character instances; the complete geometric proof above has the recorded root and independent peer audits.
The initial local stage and the symbolic-power simplification
failure are preserved under
\path{records/attempts/e6_unselected_iii_extension/}.

\end{enumerate}
\end{proof}